% Section 5.3, from lectures/L05/appendix/b_exercises.md.

\section{Exercises}
\label{c5:sec:exercises}\label{c5:app:b}

\begin{aside}
\textbf{How to work these.} Exercise~\ref{c5:ex:1} reinforces \secref{c5:sec:uart_regs}, and
Exercise~\ref{c5:ex:2} completes \code{uart_top} from \secref{c1:sec:uart_top}. The \code{fifo} the
register bank owns was built in \chapref{4}, so it is already in \file{hw/} and ready to
instantiate. Build each module from its section.

\textbf{How to check your work.} Check each module with its provided testbench, run from
\file{hw/}; the preface's \emph{Setting up} has the full flow.

The answers to every part that is not code are in \secref{sol:c5}. Try each part before you read its
answer.
\end{aside}

% ------------------------------------------------------------------------------------------
\exercise{\code{uart_regs}}{VHDL}
\label{c5:ex:1}

\xpart{a} Write your own \code{uart_regs.vhd} from the specification in \secref{c5:sec:uart_regs}
and run its testbench:

\begin{shell}
cd hw
ghdl -a --std=93 uart_def.vhd fifo.vhd uart_regs.vhd uart_regs_tb.vhd
ghdl -e --std=93 uart_regs_tb
ghdl -r --std=93 uart_regs_tb --assert-level=error
\end{shell}

\noindent \code{uart_def.vhd} (the register map) and \code{fifo.vhd} come first because
\code{uart_regs} reads the package and instantiates two FIFOs.

\xpart{b} Make \code{RX_DATA} pop the FIFO on read (advance \code{tail} whenever that index is
read). Re-run. Which check fails first, and why is it not the ``a bare \code{RX_DATA} read must not
pop'' check? Does \code{uart_top_tb} notice? Then explain, in terms of the SPI transport's abort
rule, why a read with a side effect is a genuinely harder thing to get right than the pure read plus
a separate \code{RX_POP}.

\xpart{c} \code{STATUS} is computed, not stored. Name the four things each of its bits is derived
from, and say why none of them can simply be a stored register that the datapath writes.

\xpart{d} The bank acts on every \code{reg_write} it sees as a real commit. What does it rely on the
provided SPI bridge to guarantee for that to be safe, and what would break if a half-finished
(aborted) transaction could reach it as a \code{reg_write}?

% ------------------------------------------------------------------------------------------
\exercise{\code{uart_top}}{VHDL}
\label{c5:ex:2}

\xpart{a} You built the \code{uart_top} skeleton in \chapref{1} and added \code{baud_gen},
\code{uart_tx}, \code{sync} and \code{uart_rx} across Chapters~2 and~4, declaring each block's
signals as you went. Now do it once more for \code{uart_regs}, the last block the top was waiting
for. Three signals are new here:

\begin{narrowtable}{@{}llll@{}}
  \thead{Signal} & \thead{Type} & \thead{Driven by} & \thead{Read by} \\
  \midrule
  \code{tx_empty} & \code{std_logic} & \code{uart_regs} & the feeder below \\
  \code{rx_full} & \code{std_logic} & \code{uart_regs} & nothing, in this build \\
  \code{tx_pop} & \code{std_logic} & the feeder below & \code{uart_regs} \\
\end{narrowtable}

Every other port below is a signal you already declared in Chapters~1, 2 or~4 --- which is the
point of having declared them where they were needed.

\modulefigure{uart_regs}

\begin{booktable}{@{}rp{6.4em}p{1.8em}p{15.2em}L@{}}
  \thead{\#} & \thead{\code{uart_regs} port} & \thead{Dir} & \thead{Type} & \thead{Connect to} \\
  \midrule
  1 & \code{clock} & in & \code{std_logic} & \code{clock} \\
  2 & \code{reset_s2_n} & in & \code{std_logic} & \code{reset_s2_n} \\
  3 & \code{reg_addr} & in & \code{std_logic_vector(3 downto 0)} & \code{reg_addr} (\chapref{1}) \\
  4 & \code{reg_wdata} & in & \code{std_logic_vector(31 downto 0)} & \code{reg_wdata}
    (\chapref{1}) \\
  5 & \code{reg_write} & in & \code{std_logic} & \code{reg_write} (\chapref{1}) \\
  6 & \code{tx_pop} & in & \code{std_logic} & \code{tx_pop} \\
  7 & \code{rx_byte} & in & \code{std_logic_vector(7 downto 0)} & \code{rx_byte} (\chapref{4}) \\
  8 & \code{rx_push} & in & \code{std_logic} & \code{rx_push} (\chapref{4}) \\
  9 & \code{tx_busy} & in & \code{std_logic} & \code{tx_busy} (\chapref{2}) \\
  10 & \code{frame_err} & in & \code{std_logic} & \code{frame_err} (\chapref{4}) \\
  11 & \code{reg_rdata} & out & \code{std_logic_vector(31 downto 0)} & \code{reg_rdata}
    (\chapref{1}) \\
  12 & \code{baud_div} & out & \code{std_logic_vector(15 downto 0)} & \code{baud_div}
    (\chapref{2}) \\
  13 & \code{tx_byte} & out & \code{std_logic_vector(7 downto 0)} & \code{tx_byte} (\chapref{2}) \\
  14 & \code{tx_empty} & out & \code{std_logic} & \code{tx_empty} \\
  15 & \code{rx_full} & out & \code{std_logic} & \code{rx_full} \\
\end{booktable}

\begin{vhdlcode}
-- The register bank: the bridge's bus on one side, the datapath on the other.
-- This closes every loop in the design - the bus now has something answering on
-- it, the transmitter has a byte source, and the receiver has somewhere to put
-- what it recovers.
uart_regs: entity work.uart_regs
    port map(clock, reset_s2_n, reg_addr, reg_wdata, reg_write, tx_pop, rx_byte,
             rx_push, tx_busy, frame_err, reg_rdata, baud_div, tx_byte, tx_empty,
             rx_full);
\end{vhdlcode}

\noindent Delete the \code{reg_rdata <= (others => '0');} placeholder from \chapref{1} in the same
edit: the bank drives that vector now, and leaving the placeholder in gives it two drivers.

With \code{tx_empty} finally driven, the \textbf{TX feeder} can be written, the one piece of real
logic the top has been waiting to hold. It is two concurrent statements:

\begin{vhdlcode}
-- The TX feeder: load the transmitter whenever a byte is queued and the line is
-- free, and pop the FIFO on that same edge, so exactly one byte moves per idle
-- transmitter. The instant the transmitter goes busy, tx_load drops - which is
-- what stops a second byte following it.
tx_load <= (not tx_empty) and (not tx_busy);
tx_pop  <= tx_load;
\end{vhdlcode}

\noindent Note why this could not have been written earlier: \code{tx_empty} comes from the register
bank and \code{tx_busy} from the transmitter, so the feeder needs a block from \chapref{2} and a
block from \chapref{5} to exist at the same time. That is the general shape of this design, and the
reason signals are declared where they are needed rather than all at once in \chapref{1}.

The receive side needs no equivalent: \code{uart_rx}'s \code{valid} went straight to \code{rx_push}
in \chapref{4}, so the receiver pushes its own byte into the RX FIFO.

Then run the system testbench:

\begin{shell}
cd hw
ghdl -a --std=93 uart_def.vhd reset_sync.vhd baud_gen.vhd sync.vhd fifo.vhd \
     uart_tx.vhd uart_rx.vhd uart_regs.vhd spi_slave.vhd spi_reg_bridge.vhd \
     uart_top.vhd uart_top_tb.vhd
ghdl -e --std=93 uart_top_tb
ghdl -r --std=93 uart_top_tb --assert-level=error
\end{shell}

\noindent Everything the design contains appears on that line, in dependency order, including
\code{fifo}, which \code{uart_top} never instantiates itself: \code{uart_regs} owns both of them.
With the last block in place, the system testbench elaborates the whole peripheral for the first
time.

If \code{uart_top_tb} fails on the \code{BAUD_DIV} read-back with every \code{1} bit missing from the
value you wrote (\code{got 0x0000} for the \code{0x0004} it wrote), the placeholder from \chapref{1}
is still there: each \code{'1'} the bank drives is resolving against its \code{'0'} to \code{'X'}, and
\code{to_hex} prints an \code{'X'} as a zero.

\xpart{b} The provided \code{reset_sync} asserts asynchronously but releases synchronously. Explain
what could go wrong if it released \code{reset_s2_n} asynchronously too (the moment \code{reset_n}
rises), and why asserting asynchronously is nonetheless the right choice.

\xpart{c} Why can \code{reset_sync} not just be a \code{sync} instance, given that both are two
flip-flops in series? Name the port that settles it, and describe the circular dependency you would
create by trying.

\xpart{d} Study the TX feeder you wrote in a): \code{tx_load <= (not tx_empty) and (not tx_busy)}.
Argue that it loads \textbf{exactly one} byte each time the transmitter goes idle with the FIFO
non-empty, never zero and never two. What holds \code{tx_load} low for the rest of the frame?

\xpart{e} The system testbench loops \code{tx} back to \code{rx}. Trace a single byte: an SPI write
of \code{TX_DATA}, through the FIFO and the transmitter, out on \code{tx}, back in on \code{rx},
through the receiver into the RX FIFO, and out again on an SPI read of \code{RX_DATA}. Name the
block that acts at each step.
